\documentclass[aspectratio=169,11pt]{beamer}
\usepackage[utf8]{inputenc}
\usepackage[T1]{fontenc}
\usepackage[english]{babel}
\usepackage{amsmath,amssymb}
\usepackage{graphicx}
\usepackage{booktabs}
\usetheme{Madrid}
\title[PINN et ballottement]{Improving a PINN for Violent Sloshing Simulation}
\subtitle{Interface libre, échantillonnage adaptatif et entraînement temporel}
\author{[Author(s)]}
\institute{[Laboratory / institution]}
\date{[Date]}

\begin{document}
\begin{frame}
  \titlepage
\end{frame}

\begin{frame}{Outline}
  \tableofcontents
\end{frame}

\section{Background and Objectives}
\begin{frame}{Why is violent sloshing challenging?}
\begin{itemize}
  \item Interface libre fortement déformée ; vagues raides, impacts et éventuelle rupture.
  \item Besoin de prédire vitesse, pression et fraction volumique.
  \item Main issue: a diffuse or incorrectly located PINN interface.
\end{itemize}
\vspace{0.4cm}
\textbf{Question :} quelles modifications améliorent l'interface sans utiliser les champs interFoam pendant l'entraînement ?
\end{frame}

\begin{frame}{Objectives and contributions}
\begin{itemize}
  \item Ajouter quatre mécanismes activables indépendamment.
  \item Comparer précision, robustesse et coût de calcul.
  \item Conserver une procédure reproductible et un PINN autonome.
\end{itemize}
\textit{Remplacer cette diapositive par les contributions réellement validées.}
\end{frame}

\section{Physical problem et PINN}
\begin{frame}{Physical problem}
\begin{itemize}
  \item Géométrie du réservoir : [dimensions].
  \item Excitation : [translation/rotation, amplitude, fréquence].
  \item Fluides : [densités, viscosités, tension superficielle].
  \item Conditions initiales et limites : [à préciser].
\end{itemize}
\begin{center}
\fbox{\parbox{0.75\linewidth}{Insert a tank diagram showing the initial liquid level and imposed motion.}}
\end{center}
\end{frame}

\begin{frame}{Baseline PINN}
\begin{itemize}
  \item Entrées : coordonnées normalisées $(x^*,y^*,t^*)$.
  \item Sorties : $(u^*,v^*,p^*,\alpha)$.
  \item Réseau dense Tanh ; $\alpha$ bornée par sigmoïde.
  \item Pertes : continuité, transport de $\alpha$, quantité de mouvement, initiale, parois et jauge de pression.
\end{itemize}
\[
\mathcal L=\lambda_c\mathcal L_c+\lambda_\alpha\mathcal L_\alpha+
\lambda_m\mathcal L_m+\lambda_0\mathcal L_0+\lambda_w\mathcal L_w+\lambda_p\mathcal L_p
\]
\end{frame}

\section{Proposed Methods}
\begin{frame}{Option 1 — Fourier features}
\begin{itemize}
  \item Transformer les coordonnées en composantes sinusoïdales multi-échelles.
  \item Hypothesis: better representation of steep spatial and temporal gradients.
  \item Ablation : mêmes données d'entraînement, mêmes budgets et graines.
\end{itemize}
\[
\gamma(z)=[z,\sin(2\pi Bz),\cos(2\pi Bz)]
\]
\textit{Insérer un schéma ou les fréquences réellement utilisées.}
\end{frame}

\begin{frame}{Option 2 — Adaptive sampling}
\begin{itemize}
  \item Generate a pool of candidate collocation points.
  \item Prioritize the interface ($\alpha\approx0.5$), large gradients, and high residuals.
  \item Retain a uniform fraction of points to cover the full domain.
\end{itemize}
\textbf{À surveiller :} coût de calcul des résidus et diversité des points.
\end{frame}

\begin{frame}{Option 3 — Temporal windows}
\begin{itemize}
  \item Train successively over time windows.
  \item Initialize from the previous window's weights.
  \item Use overlap or replay samples to reduce temporal forgetting.
\end{itemize}
\textit{Insérer un diagramme des fenêtres et préciser la stratégie de causalité.}
\end{frame}

\begin{frame}{Option 4 — Adaptive loss weighting}
\begin{itemize}
  \item Estimate the relative gradient scales of the different loss terms.
  \item Update weights using smoothing and bounds.
  \item Comparer aux poids fixes de la configuration de référence.
\end{itemize}
\textbf{Point de vigilance :} les poids adaptatifs doivent rester stables et interprétables.
\end{frame}

\begin{frame}{Optimization: Adam followed by L-BFGS}
\begin{itemize}
  \item Adam for the initial optimization phase.
  \item L-BFGS refinement with fixed collocation points and loss weights inside the closure.
  \item Comparer le temps total, la convergence et les résidus finaux.
\end{itemize}
\end{frame}

\section{Experimental Protocol and Results}
\begin{frame}{Outline expérimental}
\begin{enumerate}
  \item PINN de base.
  \item Chaque option activée individuellement.
  \item Combinaisons retenues et répétitions avec plusieurs graines.
  \item InterFoam utilisé comme référence d'évaluation, pas comme donnée d'entraînement.
\end{enumerate}
\end{frame}

\begin{frame}{Evaluation metrics}
\begin{itemize}
  \item Erreur de $\alpha$ et distance entre interfaces $\alpha=0.5$.
  \item Erreurs de vitesse/pression si une référence existe.
  \item Résidus physiques et stabilité pendant les impacts.
  \item Temps d'entraînement, temps par époque et mémoire utilisée.
\end{itemize}
\end{frame}

\begin{frame}{Results — summary table}
\centering
\begin{tabular}{lccc}
\toprule
Configuration & Erreur interface & Temps d'entraînement & Résidu physique \\
\midrule
PINN de base & 0.490 & 45.083 & 3724.733 \\
Fourier & 0.486 & 17.317 & 4207.534 \\
Échantillonnage adaptatif & 0.491 & 14.481 & 3876.398 \\
Fenêtres temporelles & 0.492 & 13.475 & 3641.053 \\
Pertes adaptatives & 0.492 & 22.982 & 3680.937 \\
Combinaison finale (ff=1,as=1,ts=1,alw=0) & 0.480 & 6.685 & 4169.999 \\
\bottomrule
\end{tabular}
\end{frame}

\begin{frame}{Visualizations to include}
\begin{itemize}
  \item Contours de $\alpha$ à plusieurs instants et contour $\alpha=0.5$.
  \item Comparaison PINN / interFoam sur les mêmes instants.
  \item Zoom sur l'impact et les zones de forte courbure.
  \item Courbes d'apprentissage et temps cumulés.
\end{itemize}
\end{frame}

\section{Conclusion}
\begin{frame}{Conclusions and next steps}
\begin{itemize}
  \item Main finding: [complete after measurements].
  \item Modification la plus efficace : [à compléter].
  \item Limites et travaux futurs : [à compléter].
\end{itemize}
\vfill
\centering \Large Questions ?
\end{frame}

\end{document}
